\chapter{Conclusions}
\label{cap:conclusoes}

\section{Summary of the results}
\label{sec:conc-summary}

This work studied the collinear spin correlations in QCD jets with CMS Open Data datasets, using 
two observables sensitive to the polarisation of an intermediate gluon. The first 
computed observable is the azimuthal angle $\dpsi$ between the planes of two
successive splittings of the Lund declustering. The second is the azimuthal
modulation of the three-point energy correlator. Both of them were measured 
in jets with $\pt > 700$~GeV in the 2016 collision data and compared with the corresponding
\textsc{Pythia}~8 simulation, before and after a flavour selection designed to
enrich the samples in $g \to q\bar{q}$ splittings. The results are summarised below 
in relation to the objectives of Section~\ref{sec:objetivos}.

A fully complete and reproducible analysis was developed over the released data. It
extracts the jets and their constituents information inside the CMS software environment, and
performs the declustering, the generator-level labelling of the splittings, the
computation of the energy correlator and the flavour selection outside it, with
public tools. The simulation describes the collision data kinematics, and the primary and secondary Lund jet planes
show the structure expected from the emission pattern of QCD and from the
boundaries set by the grooming condition and by the jet radius. The momentum fractions 
distributions of the two selected splittings agree between data and simulation as well, 
with mean values of $0.270$ and $0.283$ in data
against $0.268$ and $0.283$ in simulation.

The observed $\dpsi$ shows a clear $\cos(2\dpsi)$ modulation even before any flavour selection,
with $\langle \cos 2\dpsi \rangle = -0.078 \pm 0.003$ in data and
$-0.083$ in simulation. In simulation, the jets labelled as $g \to q\bar{q}$ have a
modulation about twice as strong as the background, $-0.167$ against $-0.081$,
with a negative sign as predicted for this channel. However, every channel carries a negative
modulation, including those in which little or no spin correlation is
expected. This non-spin component is largest for asymmetric splittings and
dominates the inclusive modulation, in line with the small difference between
the \textsc{Pythia}~8 predictions with and without spin correlations reported by
CMS. The three-point energy correlator also shows a modulation that the simulation describes
well in data, $+0.0200 \pm 0.0012$ against $+0.0186 \pm 0.0001$, but it has a positive sign 
and is almost independent of the channel, in contradiction with the theory prediction.
In the configuration used here,
it is therefore not a channel-selective probe of the spin correlations.

A two-stage flavour selection was developed, in which a first network separates
matched from unmatched splittings and a second network identifies the
$g \to q\bar{q}$ channel among the matched ones. This selection raises the purity 
of the selected samples from $0.03$ before any selection to between $0.48$ and $0.62$ at
the highest threshold. It performs best among the methods studied at low
efficiency, while a four-class network developed as an alternative test, similar to the one used by CMS, is better
at high efficiency. The purity counted against all jets saturates at about $0.5$,
because the tightest requirements retain proportionally more jets without a
generator-level counterpart. This sets the practical limit of the selection with
the features used.

Using this selection, the modulation of the selected jets strengthens steadily with
the threshold, from $-0.122$ at $0.30$ to $-0.266$ at $0.60$ in simulation. The data 
follows the simulation within one standard deviation up to a threshold of
$0.50$, beyond which their statistical precision is insufficient. However,
a modulation of a pure $g \to q\bar{q}$ sample could not be established. The modulation of each
labelled class changes with the threshold, so the selected samples do not behave
as mixtures of classes with fixed modulations,
and a linear extrapolation to a purity of one is not justified. The most reasonable
quantity is the difference between the $g \to q\bar{q}$ class and the background,
which stays close to $0.13$ up to a threshold of $0.50$, and which is less than 
half of the difference expected at leading order. The shift present in both classes can be explained
by a selection of configurations with a stronger non-spin
component, or by genuine $g \to q\bar{q}$ splittings that the label assigns to
other channels. Despite of these hypothesis, the samples available do not allow 
them to be studied or tested.

The sources of systematic uncertainties were identified, with those related to the purity 
limiting the interpretation most. These are the dependence of
the purity on the generator and on the definition of the channel label.
However, the most important limitation of this work is the absence of a simulated dataset 
without spin correlations processed using the selection adopted here. A dataset of that 
kind allowed CMS to establish the presence of the correlations, and it is not
available among the released simulations of the period studied. Nevertheless, the present
analysis reproduces the strategy of the CMS measurement with public data,
and its results are consistent with the CMS observation that the data agree with
\textsc{Pythia}~8 with spin correlations, although, it cannot establish by itself 
their presence. Its main contribution is the decomposition of the modulation by channel
class, which shows that the modulation of a flavour-selected sample cannot be
interpreted as that of the $g \to q\bar{q}$ channel alone. This applies to any
measurement of this kind, and it supports comparing the data with predictions
without spin correlations rather than interpreting the selected modulation
directly.

\section{Future work}
\label{sec:conc-future}

The most important extension would be a prediction without spin correlations, 
processed with the same selection as in data. In \textsc{Pythia}~8 the correlations can
be switched off in the parton shower, and \textsc{Herwig}~7 treats them explicitly
through the spin-density matrix, so that samples with and without correlations can
be generated with both. A full simulation of the CMS detector is computationally
demanding, but a first study can be carried out at particle level, where the
flavour selection, the labelling and the observables can be applied without
change. Such samples would separate the spin and non-spin components of the
modulation, and would distinguish between the two explanations of the effect of
the selection discussed in Section~\ref{sec:disc-selection}.

The flavour selection procedure could also be improved in two directions.
Although $\dpsi$ is not an input of the classifier, the shift that the selection induces in every
class shows that its score is correlated with $\dpsi$ through other features. Training another
classifier explicitly decorrelated from $\dpsi$, using, for instance, a penalty
on the distance correlation between its score and
$\dpsi$~\cite{Kasieczka:2020disco} or with an adversarial network~\cite{Louppe:2017pivot}, 
would limit this shift, possibly at the cost of
some discriminating power. Another and a simpler solution would be training the network
on a sample without spin correlations, as done by CMS. This would also ensure that any
residual correlation would not originate from spin. The dependence of the purity on the definition of the channel label
should also be studied, by comparing the matching of the two branches of the
second splitting used here with alternatives such as the matching of subjets 2 and
3 used by CMS. A second generator would finally test how much the composition and
the purity depend on the description of the parton shower and of the
hadronisation.

On the collision data side, the statistical precision could be increased by using 
other datasets of other periods of the 2016 data taking available in CMS Open Data.
This would permit a comparison with simulation at the highest thresholds. The
systematic uncertainties identified in Section~\ref{sec:systematic} should be
propagated to the results, and the analysis choices that differ from the CMS
measurement, in particular the thresholds on the transverse momentum of the
constituents and of the second subjet, should be varied. An unfolding of the
measured distributions to particle level would finally allow a direct comparison
with theoretical predictions and with the CMS result.

The three-point energy correlator requires a different treatment in order to become 
sensitive to the $g \to q\bar{q}$ channel. Its angular window was transposed from
electron-positron collisions, and its optimisation for jets at a hadron collider,
together with predictions for proton-proton collisions, would clarify the origin
of the positive modulation observed here. Restricting the triplets to the
constituents of the branches selected by the Lund declustering would also tie the
correlator to the labelled splitting, and allow the flavour selection to act on
it in the same way as on $\dpsi$.

Concluding, this analysis chain is not restricted to proton-proton collisions.
CMS also released heavy-ion data, in which the same observables 
could be used to study whether the spin correlations of the parton shower are
modified by the quark-gluon plasma. However the simulated heavy-ion datasets, that 
are suitable for this analysis, do not include jet quenching and their maximum
jet radius is at maximum of $R=0.5$. Nevertheless,
they can provide a reference without medium effects, but not a prediction of how
the medium modifies the correlations, which would require dedicated samples
generated with a model of jet quenching. The reference provided by the
proton-proton measurement, together with the understanding of the selection
obtained in this work, is a necessary starting point for such a study.